\section{Introduction}
Forecasting a stochastic dynamical system is a conditional sampling problem. Given recent observations, a useful model must represent several possible futures and the relative uncertainty between them. Generating large ensembles becomes expensive when every member requires dozens of neural evaluations.

Stochastic-interpolant forecasting starts a generative diffusion at the observed state and transforms that point mass into a forecast distribution \citep{chen2024forecasting}. Flow map matching learns finite-time deterministic transport \citep{boffi2025flowmaps}. It\^o maps extend this idea to stochastic dynamics by including the driving Brownian path as an input \citep{pan2026itomaps}. These developments suggest an efficient conditional forecast sampler. Their combination nevertheless requires care: a forecasting model starts from a fixed observation, its regression target depends on the chosen interpolant, and composition must use compatible Brownian information.

Our contribution is a conditional formulation and implementation connecting these ingredients. We derive a two-time map for the existing forecasting SDE, explain how its loss differs from a Gaussian-base construction, and give a deterministic baseline under the same map interface. We separate errors due to drift estimation, Brownian compression, and map learning. We evaluate the maps on stochastically forced Navier--Stokes trajectories. The empirical question is whether a small number of learned transitions retains forecast quality at lower inference cost. Fewer network evaluations alone do not establish this claim.
